% How accurate LIGHT is: an overview of the basic network's accuracy, its
% system map, and its reaction rates (Appendix A).
% No number here is typed by hand. They come from six generated files:
%   system_map_basic_numbers.tex   make_system_map_basic.py
%   overview_rates.tex, overview_fast.tex, overview_slow.tex
%                                  make_overview.py
%   energy_ladder_v0.tex           make_energy_ladder.py
%   REVIEW/review_numbers.tex      REVIEW/make_review.py (v0's default
%                                  network at Tahoma, the bottom line)
% Build from LIGHT/docs:
%   python3 make_system_map_basic.py && python3 make_overview.py
%   python3 make_energy_ladder.py
%   latexmk -pdf system_map_basic.tex && latexmk -pdf accuracy_overview.tex
\documentclass[11pt]{article}
\usepackage[a4paper,margin=2.2cm]{geometry}
\usepackage[T1]{fontenc}
\usepackage[scaled=0.95]{helvet}
\renewcommand{\familydefault}{\sfdefault}
\usepackage{amsmath}
\usepackage{sansmath}
\usepackage{booktabs,tabularx,array,longtable}
\usepackage{xcolor}
\usepackage{pdfpages}
\usepackage{float}
\usepackage{enumitem}
\usepackage{tikz}
\usetikzlibrary{arrows.meta,decorations.pathmorphing}
\setlist{nosep, leftmargin=1.2em}
\setlength{\parindent}{0pt}
\setlength{\parskip}{5pt}
\definecolor{ink2}{HTML}{52514E}

% LIGHT basic-network map numbers. GENERATED by make_system_map_basic.py
% from the model (0.1.0). Do not edit. Rerun the script instead.
\providecommand\lmdef[2]{\expandafter\def\csname lm@#1\endcsname{#2}}
\lmdef{e.P}{100.0}
\lmdef{e.pen}{11.2}
\lmdef{e.G1}{90.2}
\lmdef{e.C1}{11.7}
\lmdef{e.C2}{0.3}
\lmdef{e.G6}{8.9}
\lmdef{e.dU}{0.0}
\lmdef{e.cap}{0.0}
\lmdef{C1.kr}{10.0}
\lmdef{C1.heat}{1.7}
\lmdef{krs.in}{100.2}
\lmdef{krs.harp}{61.8}
\lmdef{krs.G2}{37.9}
\lmdef{krs.left}{0.5}
\lmdef{brs.harp}{61.6}
\lmdef{brs.G2}{37.9}
\lmdef{harp.B}{21.4}
\lmdef{harp.C}{13.1}
\lmdef{harp.F}{5.5}
\lmdef{harp.heat}{21.8}
\lmdef{G2.ion}{26.8}
\lmdef{C2.ion}{0.3}
\lmdef{mix.net}{2.9}
\lmdef{mix.netC}{2.9}
\lmdef{mix.heat}{0.0}
\lmdef{mix.BC}{11.5}
\lmdef{mix.CB}{14.4}
\lmdef{B.G4}{6.4}
\lmdef{B.G5}{13.8}
\lmdef{B.Q}{4.0}
\lmdef{B.QF2}{1.0}
\lmdef{B.QKR}{0.0}
\lmdef{B.QE}{3.0}
\lmdef{B.left}{0.0}
\lmdef{C.G4C}{1.1}
\lmdef{C.G5C}{3.6}
\lmdef{C.Q}{5.5}
\lmdef{C.QF2}{0.8}
\lmdef{C.QKR}{1.1}
\lmdef{C.QE}{3.6}
\lmdef{C.left}{0.0}
\lmdef{G4.hv}{5.2}
\lmdef{G4.F}{0.8}
\lmdef{G4.heat}{0.4}
\lmdef{G4C.hv}{0.8}
\lmdef{G4C.F}{0.1}
\lmdef{G4C.heat}{0.1}
\lmdef{QB.F}{0.8}
\lmdef{QB.heat}{3.2}
\lmdef{QC.F}{0.9}
\lmdef{QC.heat}{4.6}
\lmdef{tri.k2f}{11.6}
\lmdef{tri.heat}{5.8}
\lmdef{k2f.C3}{5.6}
\lmdef{k2f.Q2F}{6.0}
\lmdef{k2f.left}{0.0}
\lmdef{brk.C3}{48.0}
\lmdef{brk.Q2F}{51.9}
\lmdef{C3.hv}{4.1}
\lmdef{C3.F}{1.1}
\lmdef{C3.heat}{0.4}
\lmdef{Q2F.F}{1.2}
\lmdef{Q2F.heat}{4.8}
\lmdef{out.248}{5.2}
\lmdef{out.275}{0.8}
\lmdef{out.415}{4.1}
\lmdef{out.heat}{42.9}
\lmdef{out.elastic}{8.9}
\lmdef{out.ions}{27.1}
\lmdef{out.F}{10.4}
\lmdef{out.other}{0.5}
\lmdef{version}{0.1.0}
\lmdef{p}{1}
\lmdef{xF2}{0.5}
\lmdef{Tg0}{300}
\lmdef{P}{300}
\lmdef{pulse}{1}
\lmdef{tend}{2}
\lmdef{ne}{10^{15}}
\lmdef{N}{2.45{\times}10^{19}}
\lmdef{Edep}{0.30}
\lmdef{Te}{2.03}
\lmdef{TeMid}{2.00}
\lmdef{rTail}{0.30}
\lmdef{TtMid}{0.60}
\lmdef{TgEnd}{604}
\lmdef{pEnd}{2.03}
\lmdef{F2burn}{715}
\lmdef{F2left}{0.04}
\lmdef{krsFormed}{1.89{\times}10^{17}}
\lmdef{krfFormed}{1.05{\times}10^{17}}
\lmdef{k2fFormed}{5.30{\times}10^{16}}
\lmdef{krf.eff}{28.0}
\lmdef{brf.G4}{18.7}
\lmdef{brf.G4C}{3.2}
\lmdef{brf.tri}{50.5}
\lmdef{brf.e}{19.0}
\lmdef{brf.F2}{5.4}
\lmdef{brf.Kr}{3.2}
\lmdef{brf.QB}{11.6}
\lmdef{brf.QC}{16.1}
\lmdef{k.harp}{8.47{\times}10^{-10}}
\lmdef{pc.harpY}{90}
\lmdef{pc.harpB}{56}
\lmdef{pc.harpC}{34}
\lmdef{pc.harpX}{10}
\lmdef{pc.fC}{38}
\lmdef{k.G2}{5.0{\times}10^{-10}}
\lmdef{tauB}{6.99}
\lmdef{tauC}{50}
\lmdef{tauC3}{220}
\lmdef{gapBC}{80}
\lmdef{k.BC}{1.05{\times}10^{-11}}
\lmdef{Keq}{1.10}
\lmdef{k.G5}{1.9{\times}10^{-30}}
\lmdef{k.G5C}{3.79{\times}10^{-31}}
\lmdef{tpow.G5}{-3.32}
\lmdef{k.Q2F}{1.2{\times}10^{-10}}
\lmdef{hv.G4C}{4.5}
\lmdef{hv.G4}{5.00}
\lmdef{notes.krs}{101.9}
\lmdef{notes.krp}{27.1}
\lmdef{notes.emit}{10.1}
\lmdef{om.kr2}{8.0}
\lmdef{om.ions}{5.3}
\lmdef{om.electron}{3.8}
\lmdef{om.f2e}{5.7}
\lmdef{om.all}{8.0}
\lmdef{om.basic}{5.2}
\lmdef{om.v0}{16.8}

% LIGHT reaction-rate tables. GENERATED by make_overview.py from
% LIGHT v0 (0.1.0). Do not edit. Rerun the script instead.
\providecommand\ovdef[2]{\expandafter\def\csname ov@#1\endcsname{#2}}
\ovdef{rates.teMid}{2.00}
\ovdef{rates.rMid}{0.30}
\ovdef{rates.nBasic}{24}
\ovdef{rates.nGroups}{19}
\ovdef{rates.groups}{kr2, ions, electron, f2e}
\ovdef{rates.v0groups}{quench, kr2, ions, electron, bc, f2e}
\def\ratesBasic{G1 & \el{} + Kr \ra \el{} + \Krs{} & $\langle\sigma v\rangle$, Table A.3 & $T_e$ & Biagi v7.1, the whole 5s manifold at 9.915 eV\\
C1a & \el{} + Kr \ra \el{} + \Krs{} & $\langle\sigma v\rangle$, Table A.3 & $T_e$ & Biagi v7.1 5p group at 11.30 eV, prompt cascade to 5s\\
C1b & \el{} + Kr \ra \el{} + \Krs{} & $\langle\sigma v\rangle$, Table A.3 & $T_e$ & Biagi v7.1 d+p group at 11.998 eV, as C1a\\
C1c & \el{} + Kr \ra \el{} + \Krs{} & $\langle\sigma v\rangle$, Table A.3 & $T_e$ & Biagi v7.1 higher group at 12.75 eV, as C1a\\
C2 & \el{} + Kr \ra 2\el{} + \Krp{} & $\langle\sigma v\rangle$, Table A.3 & $T_e$ & Biagi v7.1 ionization at 13.996 eV\\
G6 & \el{} + Kr, elastic & energy loss $(2m/M)\langle\sigma_m v\varepsilon\rangle$ & $T_e$, $T_g$ & Biagi v7.1 momentum transfer\\
G2 & 2\Krs{} \ra \Krp{} + Kr + \el{} & $5.0{\times}10^{-10}$ cm$^3$s$^{-1}$ &  & Kannari 1985 Table I, Penning block\\
G3 & \Krs{} + \Ft{} \ra KrF(B) + F & $4.73{\times}10^{-10}$ cm$^3$s$^{-1}$ &  & rate from Velazco, Kolts and Setser 1976 and Ueda 1984, KrF$^*$ yield 0.90 from Velazco, Kolts and Setser 1976, B share from Yu and Setser 1990\\
G3C & \Krs{} + \Ft{} \ra KrF(C) + F & $2.9{\times}10^{-10}$ cm$^3$s$^{-1}$ &  & as G3, with Yu and Setser's $k_B/k_C$ of 1.6\\
G3X & \Krs{} + \Ft{} \ra Kr + 2F & $8.47{\times}10^{-11}$ cm$^3$s$^{-1}$ &  & as G3, the KrF$^*$-free rest, assigned to Kr + 2F\\
G4 & KrF(B) \ra Kr + F + \hnu & $1.43{\times}10^{8}$ s$^{-1}$ &  & Kannari 1985 Table I. The notes write about 4 ns\\
G4C & KrF(C) \ra Kr + F + \hnu & $2.0{\times}10^{7}$ s$^{-1}$ &  & ML joint fit. Hakuta 1990 measured 60 $\pm$ 10 ns\\
BC & KrF(B) + Kr \ra KrF(C) + Kr & $1.05{\times}10^{-11}$ cm$^3$s$^{-1}$ &  & ML joint fit. Hakuta 1990's 1.2$\times10^{-10}$ is held out\\
CB & KrF(C) + Kr \ra KrF(B) + Kr & $9.55{\times}10^{-12}$ cm$^3$s$^{-1}$ & $1/K_\mathrm{eq}(T_g)$ & detailed balance of BC, with the joint fit's $K_\mathrm{eq}$\\
G5 & KrF(B) + 2Kr \ra \KtF{} + Kr & $1.9{\times}10^{-30}$ cm$^6$s$^{-1}$ & $(T_g/300)^{-3.32}$ & ML joint fit. $T_g$ law from Shui's calculation\\
G5C & KrF(C) + 2Kr \ra \KtF{} + Kr & $3.79{\times}10^{-31}$ cm$^6$s$^{-1}$ & $(T_g/300)^{-3.32}$ & ML joint fit\\
QE & KrF(B) + \el{} \ra Kr + F + \el{} & $6.6{\times}10^{-8}$ cm$^3$s$^{-1}$ &  & Trainor and Jacob's measured product, re-read with this network's lifetime\\
QEC & KrF(C) + \el{} \ra Kr + F + \el{} & $6.6{\times}10^{-8}$ cm$^3$s$^{-1}$ &  & as QE\\
QF2 & KrF(B) + \Ft{} \ra Kr + 3F & $4.75{\times}10^{-10}$ cm$^3$s$^{-1}$ &  & ML joint fit. Xu 1993 printed 3.0 $\pm$ 0.3$\times10^{-10}$ with fast mixing\\
QF2C & KrF(C) + \Ft{} \ra Kr + 3F & $3.04{\times}10^{-10}$ cm$^3$s$^{-1}$ &  & ML joint fit\\
QKR & KrF(B) + Kr \ra 2Kr + F & $1.0{\times}10^{-14}$ cm$^3$s$^{-1}$ &  & ML joint fit, at its lower bound\\
QKRC & KrF(C) + Kr \ra 2Kr + F & $8.37{\times}10^{-13}$ cm$^3$s$^{-1}$ &  & ML joint fit\\
C3 & \KtF{} \ra 2Kr + F + \hnu & $4.5{\times}10^{6}$ s$^{-1}$ &  & Xu, Gadomski and Setser 1993, 220 $\pm$ 20 ns\\
Q2F & \KtF{} + \Ft{} \ra 2Kr + F + \Ft{} & $1.2{\times}10^{-10}$ cm$^3$s$^{-1}$ &  & Xu 1993 Table I, 1.2 $\pm$ 0.2$\times10^{-10}$\\}
\def\ratesGroups{AK8\_Kr & \Krs{} + 2Kr \ra \Krts{} + Kr & $1.7{\times}10^{-32}$ cm$^6$s$^{-1}$ & $(T_g/300)^{-0.64}$ & Ueda 1984 Table V, 1.7 $\pm$ 0.2$\times10^{-32}$. $T_g$ law from Wieme and Lenaerts 1981\\
KR2\_RAD & \Krts{} \ra 2Kr + \hnu & $3.7{\times}10^{6}$ s$^{-1}$ &  & Morikawa 1989, 270 $\pm$ 5 ns, and five other measurements\\
NK7 & \Krts{} + \Ft{} \ra KrF(B) + Kr + F & $1.05{\times}10^{-10}$ cm$^3$s$^{-1}$ &  & half of Quigley and Hughes 1978, 2.1 $\pm$ 0.2$\times10^{-10}$. The even split is Kannari's\\
K2F\_BR & \Krts{} + \Ft{} \ra \KtF{} + F & $1.05{\times}10^{-10}$ cm$^3$s$^{-1}$ &  & the other half of Quigley and Hughes 1978\\
K2F\_F & \Krts{} + F \ra KrF(B) + Kr & $3.0{\times}10^{-10}$ cm$^3$s$^{-1}$ &  & Kannari 1985, an estimate. No measurement exists\\
PEN\_KR2 & 2\Krts{} \ra \Krtp{} + 2Kr + \el{} & $5.0{\times}10^{-10}$ cm$^3$s$^{-1}$ &  & Kannari 1985 Table I\\
K15 & \Krp{} + 2Kr \ra \Krtp{} + Kr & $2.5{\times}10^{-31}$ cm$^6$s$^{-1}$ & $(T_g/300)^{-0.75}$ & Kannari 1985, Boichenko 2000\\
R11p & \el{} + \Krtp{} \ra \Krs{} + Kr & $2.24{\times}10^{-7}\,T_e^{-1/2}$ cm$^3$s$^{-1}$ & $T_e$ & Kannari 1985 Table I, on a $T_e^{-1/2}$ template\\
KSTEP & \el{} + \Krs{} \ra \Krp{} + 2\el{} & $1.56{\times}10^{-7}\,T_e^{0.71}\times$ $e^{-4.09/T_e}$ cm$^3$s$^{-1}$ & $T_e$ & Chen R9 template, an estimate\\
SUPER & \el{} + \Krs{} \ra \el{} + Kr & from G1 by detailed balance & $T_e$ & detailed balance of G1, statistical weight ratio 1/12\\
ATT & \el{} + \Ft{} \ra F\ensuremath{^-} + F & $\langle\sigma v\rangle$, Table A.3 & $T_e$ & SIGLO F$_2$ set\\
DIS\_a3 & \el{} + \Ft{} \ra \el{} + 2F & $\langle\sigma v\rangle$, Table A.3 & $T_e$ & SIGLO {\unsansmath$a\,^3\Pi_u$}, repulsive, so excitation dissociates\\
DIS\_A1 & \el{} + \Ft{} \ra \el{} + 2F & $\langle\sigma v\rangle$, Table A.3 & $T_e$ & SIGLO {\unsansmath$A\,^1\Pi_u$}, as DIS\_a3\\
VIB1 & \el{} + \Ft{} \ra \el{} + \Ft{} & $\langle\sigma v\rangle$, Table A.3 & $T_e$ & SIGLO, v = 1\\
VIB2 & \el{} + \Ft{} \ra \el{} + \Ft{} & $\langle\sigma v\rangle$, Table A.3 & $T_e$ & SIGLO, v = 2\\
VIB3 & \el{} + \Ft{} \ra \el{} + \Ft{} & $\langle\sigma v\rangle$, Table A.3 & $T_e$ & SIGLO, v = 3\\
VIB4 & \el{} + \Ft{} \ra \el{} + \Ft{} & $\langle\sigma v\rangle$, Table A.3 & $T_e$ & SIGLO, v = 4\\
K26 & \Krp{} + F\ensuremath{^-} \ra KrF(B) & $2.95{\times}10^{-6}$ cm$^3$s$^{-1}$ at 1\,atm, 300\,K, unscreened & $N$, screening & Kannari 1985, Flannery turnover with Debye screening. No measurement with F$^-$ exists\\
K108 & \Krtp{} + F\ensuremath{^-} \ra KrF(B) + Kr & $2.95{\times}10^{-6}$ cm$^3$s$^{-1}$ at 1\,atm, 300\,K, unscreened & $N$, screening & as K26\\}
\def\ratesVzero{G3 & \Krs{} + \Ft{} \ra KrF(B) + F & $8.1{\times}10^{-10}$ cm$^3$s$^{-1}$ &  & Kannari 1985 Table I, all to B. The report it credits gives 7.2$\times10^{-10}$\\
G5 & KrF(B) + 2Kr \ra \KtF{} + Kr & $6.0{\times}10^{-31}$ cm$^6$s$^{-1}$ & $(T_g/300)^{-3.32}$ & Xu, Gadomski and Setser 1993 Table II, 6.0 $\pm$ 2.0$\times10^{-31}$\\
G5\_C & KrF(C) + 2Kr \ra \KtF{} + Kr & $6.0{\times}10^{-31}$ cm$^6$s$^{-1}$ & $(T_g/300)^{-3.32}$ & as G5, applied to KrF(C)\\
AK8\_Kr & \Krs{} + 2Kr \ra \Krts{} + Kr & $2.7{\times}10^{-32}$ cm$^6$s$^{-1}$ &  & REBUILD's estimate, Kannari's Ar-partner rate given to Kr\\
KR2\_RAD & \Krts{} \ra 2Kr + \hnu & $8.0{\times}10^{7}$ s$^{-1}$ &  & Kannari 1985 Table I, 8.0$\times10^{7}$ per second\\
NK7 & \Krts{} + \Ft{} \ra KrF(B) + Kr + F & $3.0{\times}10^{-10}$ cm$^3$s$^{-1}$ &  & Kannari 1985 Table I\\
K2F\_BR & \Krts{} + \Ft{} \ra \KtF{} + F & $3.0{\times}10^{-10}$ cm$^3$s$^{-1}$ &  & Kannari 1985 Table I\\
K33 & KrF(B) + Kr \ra 2Kr + F & $4.0{\times}10^{-12}$ cm$^3$s$^{-1}$ &  & Xu, Gadomski and Setser 1993, 4.0 $\pm$ 0.5$\times10^{-12}$\\
K33\_C & KrF(C) + Kr \ra 2Kr + F & $4.0{\times}10^{-12}$ cm$^3$s$^{-1}$ &  & as K33, applied to KrF(C)\\
K32 & KrF(B) + \Ft{} \ra Kr + 3F & $3.0{\times}10^{-10}$ cm$^3$s$^{-1}$ &  & Xu, Gadomski and Setser 1993, 3.0 $\pm$ 0.3$\times10^{-10}$\\
K32\_C & KrF(C) + \Ft{} \ra Kr + 3F & $3.0{\times}10^{-10}$ cm$^3$s$^{-1}$ &  & as K32, applied to KrF(C)\\
K67 & KrF(B) + \el{} \ra Kr + F + \el{} & $4.5{\times}10^{-8}$ cm$^3$s$^{-1}$ &  & Kushner, Hanson and Schneider 1989\\
K67\_C & KrF(C) + \el{} \ra Kr + F + \el{} & $4.5{\times}10^{-8}$ cm$^3$s$^{-1}$ &  & as K67, applied to KrF(C)\\
K125 & \KtF{} + \Ft{} \ra 2Kr + F + \Ft{} & $1.5{\times}10^{-10}$ cm$^3$s$^{-1}$ &  & Kannari 1985 Table I\\
K\_2Fe & \KtF{} + \el{} \ra 2Kr + F + \el{} & $1.0{\times}10^{-7}$ cm$^3$s$^{-1}$ &  & Kannari 1985, estimated there\\
BC\_Kr & KrF(B) + Kr \ra KrF(C) + Kr & $5.0{\times}10^{-10}$ cm$^3$s$^{-1}$ &  & Kannari's B-C mixing with a Kr partner, an estimate\\
CB\_Kr & KrF(C) + Kr \ra KrF(B) + Kr & $3.41{\times}10^{-10}$ cm$^3$s$^{-1}$ &  & detailed balance of BC\_Kr at 300 K\\
URAD\_C & KrF(C) \ra Kr + F + \hnu & $1.33{\times}10^{7}$ s$^{-1}$ &  & Kannari 1985 Table I, 1.33$\times10^{7}$ per second\\}
\def\ratesXsec{G1 & Kr, 5s manifold, 9.915\,eV & $1.24{\times}10^{-17}$ & $3.44{\times}10^{-13}$ & $1.11{\times}10^{-11}$ & $6.41{\times}10^{-11}$ & $3.78{\times}10^{-10}$ & 0.10\\
C1a & Kr, 5p group, 11.3\,eV & $5.1{\times}10^{-19}$ & $5.93{\times}10^{-14}$ & $2.97{\times}10^{-12}$ & $2.11{\times}10^{-11}$ & $1.48{\times}10^{-10}$ & 0.021\\
C1b & Kr, d+p group, 11.998\,eV & $2.5{\times}10^{-19}$ & $5.84{\times}10^{-14}$ & $3.77{\times}10^{-12}$ & $3.08{\times}10^{-11}$ & $2.53{\times}10^{-10}$ & $8.7{\times}10^{-3}$\\
C1c & Kr, higher group, 12.75\,eV & $9.15{\times}10^{-21}$ & $7.73{\times}10^{-15}$ & $8.8{\times}10^{-13}$ & $1.02{\times}10^{-11}$ & $1.29{\times}10^{-10}$ & $1.42{\times}10^{-3}$\\
C2 & Kr, ionization, 13.996\,eV & $8.06{\times}10^{-21}$ & $1.57{\times}10^{-14}$ & $2.28{\times}10^{-12}$ & $2.94{\times}10^{-11}$ & $4.14{\times}10^{-10}$ & $5.84{\times}10^{-4}$\\
ATT & F$_2$, attachment & $6.58{\times}10^{-9}$ & $3.51{\times}10^{-9}$ & $2.25{\times}10^{-9}$ & $1.61{\times}10^{-9}$ & $9.75{\times}10^{-10}$ & 1\\
DIS\_a3 & F$_2$, {\unsansmath$a\,^3\Pi_u$}, 3.16\,eV & $3.01{\times}10^{-12}$ & $1.13{\times}10^{-10}$ & $3.95{\times}10^{-10}$ & $7.37{\times}10^{-10}$ & $1.35{\times}10^{-9}$ & 1\\
DIS\_A1 & F$_2$, {\unsansmath$A\,^1\Pi_u$}, 4.34\,eV & $2.06{\times}10^{-13}$ & $1.93{\times}10^{-11}$ & $8.88{\times}10^{-11}$ & $1.9{\times}10^{-10}$ & $4.02{\times}10^{-10}$ & 1\\
VIB1 & F$_2$, $v=1$, 0.111\,eV & $1.77{\times}10^{-9}$ & $1.22{\times}10^{-9}$ & $8.72{\times}10^{-10}$ & $6.61{\times}10^{-10}$ & $4.27{\times}10^{-10}$ & 1\\
VIB2 & F$_2$, $v=2$, 0.219\,eV & $2.01{\times}10^{-9}$ & $1.49{\times}10^{-9}$ & $1.09{\times}10^{-9}$ & $8.37{\times}10^{-10}$ & $5.47{\times}10^{-10}$ & 1\\
VIB3 & F$_2$, $v=3$, 0.324\,eV & $1.21{\times}10^{-9}$ & $9.75{\times}10^{-10}$ & $7.35{\times}10^{-10}$ & $5.71{\times}10^{-10}$ & $3.78{\times}10^{-10}$ & 1\\
VIB4 & F$_2$, $v=4$, 0.421\,eV & $3.83{\times}10^{-10}$ & $3.27{\times}10^{-10}$ & $2.51{\times}10^{-10}$ & $1.96{\times}10^{-10}$ & $1.31{\times}10^{-10}$ & 1\\
\midrule
ATT & energy an attached \el{} takes, eV & 0.37 & 0.46 & 0.53 & 0.59 & 0.72 & \\
 & $\tfrac32 T_e$, for comparison & 0.75 & 1.50 & 2.25 & 3.00 & 4.50 & \\}
\def\ratesTe{0.5 & 1 & 1.5 & 2 & 3}
\def\speciesTable{Kr & 0 & 1.5 & ground state\\
\Ft{} & 0 & 2.5 & ground state\\
F & 0.8 & 1.5 & half the F$_2$ bond, $D_0$ = 1.60 eV\\
\Krs{} & 9.915 & 1.5 & Biagi's 5s threshold, the 5s[3/2]2 level\\
\Krp{} & 13.996 & 1.5 & Biagi's ionization threshold\\
KrF(B) & 6.15 & 2.5 & F + the 248-nm quantum + 0.35 eV of X-state repulsion\\
KrF(C) & 6.1401 & 2.5 & 80 cm$^{-1}$ below B\\
\KtF{} & 4.1 & 3 & F + about 3.0 eV of 415-nm band + about 0.3 eV\\
\Krts{} & 9.2 & 2.5 & the 147-nm excimer, an estimated well\\
\Krtp{} & 12.846 & 2.5 & IP(Kr) less $D_0$(Kr$_2^+$) = 1.15 eV\\
F\ensuremath{^-} & \ensuremath{-}2.6 & 1.5 & $E_\mathrm{form}$(F) less the electron affinity, 3.40 eV\\}

% LIGHT accuracy numbers (fast). GENERATED by make_overview.py
% from LIGHT v0 (0.1.0) and ML/. Do not edit. Rerun the script instead.
\providecommand\ovdef[2]{\expandafter\def\csname ov@#1\endcsname{#2}}
\ovdef{xu.50}{0.91}
\ovdef{xu.50.v0}{2.89}
\ovdef{xu.100}{0.97}
\ovdef{xu.100.v0}{2.41}
\ovdef{xu.50.meas}{$5.0{\times}10^{7}$}
\ovdef{xu.100.meas}{$6.2{\times}10^{7}$}
\ovdef{eq5.dev}{$1.0{\times}10^{-4}$}
\ovdef{closure.max}{$10^{-12}$}
\ovdef{closure.runs}{9}
\ovdef{sc.eden_kr}{2.3}
\ovdef{sc.eden_f2}{3.2}
\ovdef{sc.ksu_fig2}{4.2}
\ovdef{sc.ksu_fig1}{6.1}
\ovdef{sc.xu_figs5_7}{11.7}
\ovdef{sc.hakuta_fig3}{34.1}
\ovdef{sc.hakuta_fig2}{19.3}
\ovdef{sc.ys_ir300}{\ensuremath{-}0.5}
\ovdef{sc.ys_ratio_230_300}{\ensuremath{-}13.6}
\ovdef{sc.bs_tau_b}{\ensuremath{-}22.3}
\ovdef{ys.meas}{0.16}
\ovdef{ys.err}{0.02}
\ovdef{ysr.meas}{1.30}
\ovdef{ysr.err}{0.06}
\ovdef{sc.train.lo}{2.3}
\ovdef{sc.train.hi}{6.1}
\ovdef{sc.bs.meas}{9.0}
\ovdef{sc.bs.err}{0.6}
\ovdef{tah.g0.meas}{6.39}
\ovdef{tah.g0}{6.11}
\ovdef{tah.g0.x}{0.96}
\ovdef{tah.isat.meas}{1.4}
\ovdef{tah.isat}{1.54}
\ovdef{tah.isat.err}{0.1}
\ovdef{tah.isat.x}{1.10}
\ovdef{tah.eta.meas}{0.27}
\ovdef{tah.eta}{0.284}
\ovdef{tah.eta.x}{1.05}
\ovdef{tah.form}{0.318}
\ovdef{tah.pin}{330}
\ovdef{tah.torr}{665}
\ovdef{tah.f2torr}{2.8}
\ovdef{pure.Kr.x}{1.02}
\ovdef{pure.Ar.x}{1.03}
\ovdef{pure.dr.lo}{0.87}
\ovdef{pure.dr.hi}{1.08}
\ovdef{pure.dr.te.lo}{0.026}
\ovdef{pure.dr.te.hi}{2}
\ovdef{pure.k15}{$2.5{\times}10^{-31}$}
\ovdef{pure.conv.lo}{$2.74{\times}10^{-31}$}
\ovdef{pure.conv.hi}{$3.3{\times}10^{-31}$}
\ovdef{pure.vuv}{0.57}
\ovdef{pure.vuv.meas}{0.42}
\ovdef{pure.vuv.err}{0.05}
\ovdef{pure.vuv.x}{1.36}
\ovdef{pure.vuv.ideal}{0.54}
\ovdef{pure.vuv.nothermal}{0.55}
\ovdef{pure.thermal}{0.060}

% LIGHT accuracy numbers (slow). GENERATED by make_overview.py
% from LIGHT v0 (0.1.0) and ML/. Do not edit. Rerun the script instead.
\providecommand\ovdef[2]{\expandafter\def\csname ov@#1\endcsname{#2}}
\ovdef{fit.k_BC.lo}{$4.63{\times}10^{-12}$}
\ovdef{fit.k_BC.mid}{$9.8{\times}10^{-12}$}
\ovdef{fit.k_BC.hi}{$1.73{\times}10^{-11}$}
\ovdef{fit.K_eq.lo}{0.84}
\ovdef{fit.K_eq.mid}{1.07}
\ovdef{fit.K_eq.hi}{1.38}
\ovdef{fit.tau_C.lo}{41}
\ovdef{fit.tau_C.mid}{50}
\ovdef{fit.tau_C.hi}{63}
\ovdef{fit.tau_B.lo}{6.79}
\ovdef{fit.tau_B.mid}{7.02}
\ovdef{fit.tau_B.hi}{7.22}
\ovdef{fit.k_F2.lo}{$3.92{\times}10^{-10}$}
\ovdef{fit.k_F2.mid}{$4.57{\times}10^{-10}$}
\ovdef{fit.k_F2.hi}{$5.3{\times}10^{-10}$}
\ovdef{fit.k_F2_C.lo}{$2.74{\times}10^{-10}$}
\ovdef{fit.k_F2_C.mid}{$3.02{\times}10^{-10}$}
\ovdef{fit.k_F2_C.hi}{$3.32{\times}10^{-10}$}
\ovdef{fit.k3_Kr.lo}{$1.26{\times}10^{-30}$}
\ovdef{fit.k3_Kr.mid}{$1.92{\times}10^{-30}$}
\ovdef{fit.k3_Kr.hi}{$2.44{\times}10^{-30}$}
\ovdef{fit.k3_Kr_C.lo}{$1.93{\times}10^{-31}$}
\ovdef{fit.k3_Kr_C.mid}{$3.64{\times}10^{-31}$}
\ovdef{fit.k3_Kr_C.hi}{$5.12{\times}10^{-31}$}
\ovdef{fit.bshare.lo}{0.57}
\ovdef{fit.bshare.mid}{0.69}
\ovdef{fit.bshare.hi}{0.80}
\ovdef{fit.bshare.light}{0.68}
\ovdef{fit.names}{eden\_kr, eden\_f2, ksu\_fig2, ksu\_fig1, ys\_ir300}
\ovdef{post.g0.lo}{5.42}
\ovdef{post.g0.mid}{5.83}
\ovdef{post.g0.hi}{6.63}
\ovdef{post.n}{40}
\ovdef{fit.n}{10}
\ovdef{link.base}{0.284}
\ovdef{link.yield.lo}{0.256}
\ovdef{link.yield.lo.d}{\ensuremath{-}0.028}
\ovdef{link.yield.hi}{0.298}
\ovdef{link.yield.hi.d}{+0.014}
\ovdef{link.yield.arf}{0.227}
\ovdef{link.yield.arf.d}{\ensuremath{-}0.057}
\ovdef{link.nascentC}{0.307}
\ovdef{link.nascentC.d}{+0.023}
\ovdef{link.ionY}{0.241}
\ovdef{link.ionY.d}{\ensuremath{-}0.043}
\ovdef{link.ionC}{0.266}
\ovdef{link.ionC.d}{\ensuremath{-}0.018}
\ovdef{link.ionRate}{0.284}
\ovdef{link.ionRate.d}{\ensuremath{-}0.001}
\ovdef{link.kr2F}{0.268}
\ovdef{link.kr2F.d}{\ensuremath{-}0.017}
\ovdef{link.kr2split.lo}{0.281}
\ovdef{link.kr2split.lo.d}{\ensuremath{-}0.003}
\ovdef{link.kr2split.hi}{0.287}
\ovdef{link.kr2split.hi.d}{+0.003}
\ovdef{link.k3}{0.290}
\ovdef{link.k3.d}{+0.006}
\ovdef{link.minor}{0.006}
\ovdef{el.loop.te}{1.61}
\ovdef{el.loop.ne}{$6.52{\times}10^{14}$}
\ovdef{el.loop.g0}{6.05}
\ovdef{el.loop.form}{0.317}
\ovdef{el.light.te}{1.59}
\ovdef{el.light.ne}{$6.37{\times}10^{14}$}
\ovdef{el.light.g0}{6.11}
\ovdef{el.light.form}{0.318}
\ovdef{el.dev.te}{0.7}
\ovdef{el.dev.ne}{2.3}
\ovdef{el.dev.g0}{1.0}
\ovdef{el.dev.formation}{0.5}
\ovdef{el.eatt}{0.55}
\ovdef{el.eatt.te}{1.60}
\ovdef{el.eatt.32}{2.40}
\ovdef{sw.krw.lo}{0.96}
\ovdef{sw.krw.hi}{1.03}
\ovdef{sw.slow}{0.92}
\ovdef{sw.slow.err}{0.04}
\ovdef{sw.slow.sigma}{1.8}
\ovdef{sw.f2.lo}{0.48}
\ovdef{sw.f2.hi}{2.41}
\def\linksTable{KrF$^*$ yield of \Krs{} + \Ft & 0.90 & 0.70 & the measured 0.90 less its stated 0.20 & 0.256 & \ensuremath{-}0.028\\
 &  & 1.0 & the published assignment & 0.298 & +0.014\\
 &  & 0.5 & as for Ar$^*$ + F$_2$ & 0.227 & \ensuremath{-}0.057\\
share of that KrF formed in C & 0.38 & 0 & if nascent C relaxes to B at 1 atm & 0.307 & +0.023\\
KrF$^*$ yield of ion-ion recombination & 1 & 0.64 & as for Cl$^-$ & 0.241 & \ensuremath{-}0.043\\
share of that KrF formed in C & 0 & 0.38 & as for SF$_6^-$ & 0.266 & \ensuremath{-}0.018\\
ion-ion recombination rate & Flannery & \x{0.38} & Avco's model input & 0.284 & \ensuremath{-}0.001\\
\Krts{} + F \ra KrF(B) + Kr & $3.0{\times}10^{-10}$ & 0 & no measurement exists & 0.268 & \ensuremath{-}0.017\\}

% LIGHT v0 energy ladder. GENERATED by make_energy_ladder.py
% from LIGHT v0 (0.1.0). Do not edit. Rerun the script instead.
\providecommand\laddef[2]{\expandafter\def\csname lad@#1\endcsname{#2}}
\laddef{version}{0.1.0}
\laddef{E.B}{6.15}
\laddef{DeB}{5.25}
\laddef{BoverKrF}{5.35}
\laddef{wl.G4}{248}
\laddef{hv.G4}{5.00}
\laddef{heat.G4}{0.35}
\laddef{heat.G3}{2.97}
\laddef{heat.K26}{5.25}
\laddef{heat.G5}{2.05}
\laddef{wl.C3}{\ensuremath{\sim}415}
\laddef{hv.C3}{3.00}
\laddef{heat.C3}{0.30}
\laddef{gapBC}{80}
\laddef{notdrawn}{G6, VIB1, VIB2, VIB3 and VIB4}
\laddef{toe}{G2, PEN\_KR2 and SUPER}
\laddef{zero}{}
\laddef{groups}{\texttt{kr2}, \texttt{ions}, \texttt{electron} and \texttt{f2e}}
\laddef{pale}{Kr\ensuremath{_2^*}, Kr\ensuremath{_2^+} and F\ensuremath{^-}}
\laddef{wl.V}{147}
\laddef{hv.V}{8.43}
\laddef{heat.V}{0.77}
\def\ladpicture{%
\definecolor{ladE}{HTML}{2A78D6}% electron gas
\definecolor{ladX}{HTML}{EDA100}% Kr*
\definecolor{ladB}{HTML}{1BAF7A}% KrF(B), 248 and 272 nm
\definecolor{ladC}{HTML}{008300}% KrF(C), 275 nm
\definecolor{ladT}{HTML}{4A3AA7}% Kr2F*, 420 nm
\definecolor{ladV}{HTML}{A63A83}% Kr2*, 147 nm
\definecolor{ladH}{HTML}{E34948}% heat
\definecolor{ladI}{HTML}{52514E}% ions
\definecolor{ladF}{HTML}{A8A69E}% F atoms
\definecolor{ladInk}{HTML}{0B0B0B}%
\definecolor{ladInk2}{HTML}{52514E}%
\definecolor{ladMuted}{HTML}{898781}%
\tikzset{
  lad grid/.style={draw=ladMuted!22, line width=0.3pt},
  lad axisline/.style={draw=ladMuted, line width=0.5pt},
  lad lev/.style={line width=1.4pt},
  lad pair/.style={line width=0.8pt, dash pattern=on 3pt off 2pt},
  lad tick/.style={line width=0.5pt},
  lad arr/.style={line width=0.7pt, -{Triangle[width=3.2pt,length=3.4pt]}},
  lad hv/.style={line width=1.3pt, -{Triangle[width=4.2pt,length=4.4pt]}},
  lad on/.style={},
  lad opt/.style={dash pattern=on 2.4pt off 1.6pt},
  lad zero/.style={dash pattern=on 0.6pt off 1.4pt, line cap=round},
  lad heat/.style={draw=ladH, line width=0.6pt, decorate,
    decoration={snake, amplitude=0.5mm, segment length=1.4mm,
    post length=0.5mm}},
  lad lab/.style={font=\scriptsize, text=ladInk, inner sep=1pt},
  lad lab2/.style={font=\scriptsize, text=ladInk, inner sep=1pt,
    align=right},
  lad lab2l/.style={font=\scriptsize, text=ladInk, inner sep=1pt,
    align=left},
  lad alab/.style={font=\scriptsize, text=ladInk, inner sep=0.8pt,
    fill=white, fill opacity=0.85, text opacity=1, align=center},
  lad slab/.style={font=\tiny, text=ladInk, inner sep=0.6pt,
    fill=white, fill opacity=0.85, text opacity=1, align=right},
  lad hlab/.style={font=\tiny, text=ladInk2, inner sep=0.8pt,
    fill=white, fill opacity=0.85, text opacity=1},
  lad tlab/.style={font=\tiny, text=ladInk2, inner sep=0.8pt},
  lad axis/.style={font=\scriptsize, text=ladInk2, inner sep=1pt},
  lad head/.style={font=\footnotesize\bfseries, text=ladInk2},
}
\begin{tikzpicture}
\draw[lad grid] (0.950,1.000) -- (16.200,1.000);
\draw[lad grid] (0.950,3.000) -- (16.200,3.000);
\draw[lad grid] (0.950,5.000) -- (16.200,5.000);
\draw[lad grid] (0.950,7.000) -- (16.200,7.000);
\draw[lad grid] (0.950,9.000) -- (16.200,9.000);
\draw[lad grid] (0.950,11.000) -- (16.200,11.000);
\draw[lad grid] (0.950,13.000) -- (16.200,13.000);
\draw[lad grid] (0.950,15.000) -- (16.200,15.000);
\draw[lad grid] (0.950,17.000) -- (16.200,17.000);
\draw[lad axisline] (0.950,0.000) -- (0.950,17.450);
\draw[lad lev,draw=ladInk2] (2.650,3.000) -- (6.300,3.000);
\draw[lad lev,draw=ladX] (2.650,12.915) -- (6.300,12.915);
\draw[lad lev,draw=ladI] (2.650,16.996) -- (6.300,16.996);
\draw[lad lev,draw=ladV!40] (5.000,12.200) -- (6.300,12.200);
\draw[lad lev,draw=ladI!40] (5.000,15.846) -- (6.300,15.846);
\draw[lad tick,draw=ladMuted] (3.550,14.300) -- (4.400,14.300);
\draw[lad tick,draw=ladMuted] (3.730,14.998) -- (4.400,14.998);
\draw[lad tick,draw=ladMuted] (3.910,15.750) -- (4.400,15.750);
\draw[lad tick,draw=ladMuted] (5.740,3.770) -- (6.020,3.770);
\draw[lad pair,draw=ladInk2!40] (10.200,14.396) -- (12.300,14.396);
\draw[lad pair,draw=ladInk2!40] (10.200,13.246) -- (12.300,13.246);
\draw[lad pair,draw=ladInk2] (7.000,3.800) -- (12.300,3.800);
\draw[lad lev,draw=ladB] (8.400,9.180) -- (11.050,9.180);
\draw[lad lev,draw=ladC] (8.400,9.110) -- (11.050,9.110);
\draw[lad lev,draw=ladT] (7.300,7.100) -- (8.650,7.100);
\draw[lad tick,draw=ladMuted] (9.210,4.150) -- (9.490,4.150);
\draw[lad tick,draw=ladMuted] (10.010,4.640) -- (10.290,4.640);
\draw[lad tick,draw=ladMuted] (7.840,4.100) -- (8.120,4.100);
\draw[lad lev,draw=ladInk2] (13.900,3.000) -- (15.100,3.000);
\draw[lad lev,draw=ladF] (13.900,3.800) -- (15.100,3.800);
\draw[lad pair,draw=ladInk2] (13.900,4.600) -- (15.100,4.600);
\draw[lad pair,draw=ladInk2!40] (13.900,1.200) -- (15.100,1.200);
\draw[lad lev,draw=ladI!40] (13.900,0.400) -- (15.100,0.400);
\draw[lad tick,draw=ladMuted] (14.290,6.160) -- (15.100,6.160);
\draw[lad tick,draw=ladMuted] (14.570,7.340) -- (15.100,7.340);
\draw[lad arr,lad on,draw=ladE] (2.800,3.000) -- node[pos=0.52,rotate=90,lad alab] {C2} (2.800,16.996);
\draw[lad arr,lad on,draw=ladE] (3.080,3.000) -- node[pos=0.40,rotate=90,lad alab] {G1} (3.080,12.915);
\draw[lad arr,lad opt,draw=ladE] (3.360,12.915) -- node[pos=0.78,rotate=90,lad alab] {SUPER} (3.360,3.000);
\draw[lad arr,lad on,draw=ladX] (3.080,12.915) -- node[pos=0.36,rotate=90,lad alab] {G2} (3.080,16.996);
\draw[lad arr,lad opt,draw=ladE] (3.360,12.915) -- node[pos=0.62,rotate=90,lad alab] {KSTEP} (3.360,16.996);
\draw[lad arr,lad on,draw=ladE] (3.620,3.000) -- (3.620,14.300);
\draw[lad arr,lad on,draw=ladE] (3.800,3.000) -- node[pos=0.55,rotate=90,lad alab] {C1a--c} (3.800,14.998);
\draw[lad arr,lad on,draw=ladE] (3.980,3.000) -- (3.980,15.750);
\draw[lad heat] (4.400,15.750) -- (4.400,12.915);
\draw[lad arr,lad on,draw=ladH] (2.700,12.915) -- (2.350,12.065);
\draw[lad arr,lad opt,draw=ladX] (5.100,12.915) -- (5.100,12.200);
\draw[lad arr,lad opt,draw=ladI] (5.250,16.996) -- node[pos=0.50,rotate=90,lad alab] {K15} (5.250,15.846);
\draw[lad arr,lad opt,draw=ladI] (5.800,15.846) -- node[pos=0.45,rotate=90,lad alab] {R11p} (5.800,12.915);
\draw[lad arr,lad opt,draw=ladV] (6.120,12.200) -- node[pos=0.62,rotate=90,lad alab] {PEN\_KR2} (6.120,15.846);
\draw[lad hv,lad opt,draw=ladV] (5.880,12.200) -- node[pos=0.50,rotate=90,lad alab] {KR2\_RAD, 147 nm} (5.880,3.770);
\draw[lad heat] (5.880,3.770) -- (5.880,3.000);
\draw[lad arr,lad opt,draw=ladV] (6.300,12.200) -- node[pos=0.42,sloped,lad alab] {K2F\_BR} (7.620,7.100);
\draw[lad arr,lad opt,draw=ladV] (6.300,12.200) -- node[pos=0.42,sloped,lad alab] {NK7} (8.660,9.150);
\draw[lad arr,lad opt,draw=ladV] (6.300,12.200) -- node[pos=0.62,sloped,lad alab] {K2F\_F} (8.900,9.150);
\draw[lad arr,lad on,draw=ladX] (6.300,12.915) -- node[pos=0.40,sloped,lad alab] {G3, G3C} (9.140,9.150);
\draw[lad arr,lad opt,draw=ladI] (10.450,13.246) -- node[pos=0.50,rotate=90,lad alab] {K108} (10.450,9.150);
\draw[lad arr,lad opt,draw=ladI] (10.720,14.396) -- node[pos=0.62,rotate=90,lad alab] {K26} (10.720,9.150);
\draw[lad arr,lad on,draw=ladB] (8.520,9.150) -- (8.520,7.100);
\draw[lad hv,lad on,draw=ladB] (9.350,9.150) -- node[pos=0.50,rotate=90,lad alab] {G4, 248 nm} (9.350,4.150);
\draw[lad heat] (9.350,4.150) -- (9.350,3.800);
\draw[lad hv,lad on,draw=ladC] (10.150,9.140) -- node[pos=0.50,rotate=90,lad alab] {G4C, \ensuremath{\sim}275 nm} (10.150,4.640);
\draw[lad heat] (10.150,4.640) -- (10.150,3.800);
\draw[lad arr,lad on,draw=ladH] (10.930,9.140) -- (10.930,8.340);
\draw[lad hv,lad on,draw=ladT] (7.980,7.100) -- node[pos=0.64,rotate=90,lad alab] {C3, \ensuremath{\sim}415 nm} (7.980,4.100);
\draw[lad heat] (7.980,4.100) -- (7.980,3.800);
\draw[lad arr,lad on,draw=ladH] (7.520,7.100) -- (7.520,6.300);
\draw[lad arr,lad opt,draw=ladE] (14.360,3.000) -- node[pos=0.62,rotate=90,lad alab] {DIS\_a3} (14.360,6.160);
\draw[lad arr,lad opt,draw=ladE] (14.640,3.000) -- node[pos=0.66,rotate=90,lad alab] {DIS\_A1} (14.640,7.340);
\draw[lad heat] (15.100,7.340) -- (15.100,4.600);
\draw[lad arr,lad opt,draw=ladE] (14.080,3.000) -- node[pos=0.50,rotate=90,lad alab] {ATT} (14.080,1.200);
\node[lad axis,anchor=east] at (0.870,1.000) {\ensuremath{-}2};
\node[lad axis,anchor=east] at (0.870,3.000) {0};
\node[lad axis,anchor=east] at (0.870,5.000) {2};
\node[lad axis,anchor=east] at (0.870,7.000) {4};
\node[lad axis,anchor=east] at (0.870,9.000) {6};
\node[lad axis,anchor=east] at (0.870,11.000) {8};
\node[lad axis,anchor=east] at (0.870,13.000) {10};
\node[lad axis,anchor=east] at (0.870,15.000) {12};
\node[lad axis,anchor=east] at (0.870,17.000) {14};
\node[lad axis,anchor=east] at (0.870,18.000) {eV};
\node[lad head,anchor=center] at (4.475,18.000) {Kr and Kr\ensuremath{_2}};
\node[lad head,anchor=center] at (9.650,18.000) {KrF and Kr\ensuremath{_2}F};
\node[lad head,anchor=center] at (14.500,18.000) {F\ensuremath{_2}};
\node[lad lab,anchor=east] at (2.570,3.000) {Kr};
\node[lad lab,anchor=east] at (2.570,12.915) {Kr\ensuremath{^*} 9.915};
\node[lad lab,anchor=east] at (2.570,16.996) {Kr\ensuremath{^+} 13.996};
\node[lad lab2l,anchor=north west] at (5.030,12.150) {Kr\ensuremath{_2^*}\\9.2};
\node[lad lab2,anchor=south east] at (6.300,15.896) {Kr\ensuremath{_2^+}\\12.846};
\node[lad tlab,anchor=west] at (4.460,14.300) {11.30};
\node[lad tlab,anchor=west] at (4.460,14.998) {11.998};
\node[lad tlab,anchor=west] at (4.460,15.750) {12.75};
\node[lad alab,anchor=east] at (2.310,12.195) {G3X};
\node[lad alab,anchor=east] at (5.040,12.557) {AK8\_Kr};
\node[lad hlab,anchor=west] at (5.980,3.370) {0.77};
\node[lad lab,anchor=west] at (12.380,14.396) {Kr\ensuremath{^+} + F\ensuremath{^-} 11.396};
\node[lad lab,anchor=west] at (12.380,13.246) {Kr\ensuremath{_2^+} + F\ensuremath{^-} 10.246};
\node[lad lab,anchor=south east] at (12.300,3.840) {Kr + F 0.8};
\node[lad lab2l,anchor=west] at (11.130,9.150) {KrF(B) 6.15\\KrF(C) 80\,cm\ensuremath{^{-1}} below\\BC, CB mix them};
\node[lad lab2,anchor=east] at (7.220,7.100) {Kr\ensuremath{_2}F\ensuremath{^*}\\4.1};
\node[lad slab,anchor=east] at (8.470,8.125) {G5,\\G5C};
\node[lad hlab,anchor=north] at (9.350,3.420) {X, 0.35};
\node[lad hlab,anchor=north] at (10.150,3.420) {A, 0.84};
\node[lad alab,anchor=north west] at (10.620,8.290) {QF2, QKR, QE\\QF2C, QKRC, QEC};
\node[lad hlab,anchor=north] at (7.980,3.720) {0.30};
\node[lad alab,anchor=north] at (7.520,6.250) {Q2F};
\node[lad lab,anchor=east] at (13.820,3.000) {F\ensuremath{_2}};
\node[lad lab,anchor=east] at (13.820,3.800) {F 0.8};
\node[lad lab,anchor=east] at (13.820,4.600) {F + F 1.6};
\node[lad lab,anchor=east] at (13.820,1.200) {F\ensuremath{^-} + F \ensuremath{-}1.8};
\node[lad lab,anchor=east] at (13.820,0.400) {F\ensuremath{^-} \ensuremath{-}2.6};
\node[lad tlab,anchor=west] at (15.150,6.160) {{\unsansmath\ensuremath{a\,^3\Pi_u}} 3.16};
\node[lad tlab,anchor=west] at (15.150,7.340) {{\unsansmath\ensuremath{A\,^1\Pi_u}} 4.34};
\draw[lad arr,lad on,draw=ladInk2] (1.250,-0.550) -- (1.950,-0.550);
\node[lad lab,anchor=west] at (2.050,-0.550) {a row of the default network};
\draw[lad arr,lad opt,draw=ladInk2] (8.850,-0.550) -- (9.550,-0.550);
\node[lad lab,anchor=west] at (9.650,-0.550) {a row of an optional group, off by default};
\draw[lad hv,lad on,draw=ladB] (1.250,-0.970) -- (1.950,-0.970);
\node[lad lab,anchor=west] at (2.050,-0.970) {a photon, with its row and band};
\draw[lad arr,lad on,draw=ladH] (8.850,-0.970) -- (9.550,-0.970);
\node[lad lab,anchor=west] at (9.650,-0.970) {a quench, the level's energy to heat};
\draw[lad heat] (1.250,-1.390) -- (1.950,-1.390);
\node[lad lab,anchor=west] at (2.050,-1.390) {heat as a lower state falls apart (eV) or cascades};
\draw[lad pair,draw=ladInk2] (8.850,-1.390) -- (9.550,-1.390);
\node[lad lab,anchor=west] at (9.650,-1.390) {a pair, the sum of two formation energies};
\draw[lad tick,draw=ladMuted] (1.250,-1.810) -- (1.950,-1.810);
\node[lad lab,anchor=west] at (2.050,-1.810) {a threshold, or where a photon leaves its lower state};
\draw[lad lev,draw=ladInk2!40] (8.850,-1.810) -- (9.550,-1.810);
\node[lad lab,anchor=west] at (9.650,-1.810) {a species that only an optional group makes};
\end{tikzpicture}}

\input{REVIEW/review_numbers}
\makeatletter
\newcommand\lm[1]{\ifcsname lm@#1\endcsname\csname lm@#1\endcsname
  \else\errmessage{no map number '#1'. Rerun make_system_map_basic.py}\fi}
\newcommand\ov[1]{\ifcsname ov@#1\endcsname\csname ov@#1\endcsname
  \else\errmessage{no overview number '#1'. Rerun make_overview.py}\fi}
\newcommand\lad[1]{\ifcsname lad@#1\endcsname\csname lad@#1\endcsname
  \else\errmessage{no ladder number '#1'. Rerun make_energy_ladder.py}\fi}
\newcommand\rv[1]{\ifcsname rv@#1\endcsname\csname rv@#1\endcsname
  \else\errmessage{no review number '#1'. Rerun REVIEW/make_review.py}\fi}
\makeatother

% chemistry and units
\newcommand{\unit}[1]{\ensuremath{\,\mathrm{#1}}}
\newcommand{\el}{\ensuremath{\mathrm{e}^-}}
\newcommand{\Krs}{Kr\ensuremath{^*}}
\newcommand{\Krp}{Kr\ensuremath{^+}}
\newcommand{\Krtp}{Kr\ensuremath{_2^+}}
\newcommand{\Krts}{Kr\ensuremath{_2^*}}
\newcommand{\KtF}{Kr\ensuremath{_2}F\ensuremath{^*}}
\newcommand{\Ft}{F\ensuremath{_2}}
\newcommand{\Fm}{F\ensuremath{^-}}
\newcommand{\ra}{\ensuremath{\to}\ }
\newcommand{\hnu}{\ensuremath{h\nu}}
\newcommand{\Edep}{\ensuremath{E_\mathrm{dep}}}
\newcommand{\gIP}{\ensuremath{g_0I_\mathrm{sat}/P_\mathrm{in}}}
\newcommand{\x}[1]{\ensuremath{\times}#1}
% a factor of the measurement as a percentage of it (0.89 -> 89), to the
% factor's own precision: \pctof{x} for two decimals, \pctof[1]{x} for three
\newcommand{\pctof}[2][0]{\pgfmathparse{100*#2}%
  \pgfmathprintnumber[fixed,precision=#1]{\pgfmathresult}}

\begin{document}
\sansmath

{\LARGE\bfseries How accurate LIGHT is}\par
\vspace{4pt}
{\large The basic network of LIGHT v0, version \lm{version}, on 28 September
2026, with its energy ladder, its system map and its reaction rates}\par
\vspace{10pt}

\noindent\fbox{\parbox{\dimexpr\linewidth-2\fboxsep-2\fboxrule\relax}{%
\setlength{\parskip}{5pt}%
\textbf{Bottom line.} LIGHT v0 reproduces the fitted KrF decays [1, 2] within \ov{sc.train.lo}
to \ov{sc.train.hi}\,\% rms and predicts decays that no fit used [4] within \ov{sc.xu_figs5_7}\,\% rms.
At the Tahoma laser [5] its default network, which holds $n_e$ fixed and
has no charge chemistry, gives \pctof{\rv{tah.v0.def.x}}\,\% of the
measured small-signal gain of \ov{tah.g0.meas}\unit{\%/cm}, and with its
optional ion, \Ft{}-electron and \Krts{} groups on and $n_e$ free, as ML
scores it, \pctof{\ov{tah.g0.x}}\,\%.

That configuration also gives \pctof{\ov{tah.isat.x}}\,\% of Tahoma's
saturation intensity and \pctof{\ov{tah.eta.x}}\,\% of its formation
efficiency, and the Tahoma scores in Table~\ref{tab:score} use it. v0\_1
and v0\_2 carry the charge chemistry in their default networks. The C/B emission ratio of Yu and Setser [3] is matched to \ov{sc.ys_ir300}\,\%. No fit used Tahoma's numbers and no rate was set by a laser
measurement, so the Tahoma scores are predictions.

Apart from the default network's gain at Tahoma, its largest misses are on data that conflict with the fitted sets,
\ov{sc.hakuta_fig2} and \ov{sc.hakuta_fig3}\,\% rms on the C-A decays of
Hakuta et al.\ [6] and \ov{sc.bs_tau_b}\,\% on the KrF(B) lifetime of
Burnham and Searles [7], and the VUV yield of pure krypton at low pump,
\pctof{\ov{pure.vuv.x}}\,\% of Morozov's measurement [8].

{\footnotesize [1] Eden, Waynant, Searles and Burnham, Appl.\ Phys.\
Lett.\ 32, 733 (1978). [2] Gadomski, Xu, Setser and Nelson, Chem.\ Phys.\
Lett.\ 189, 153 (1992). [3] Yu and Setser, J.\ Phys.\ Chem.\ 94, 2934
(1990). [4] Xu, Gadomski and Setser, J.\ Chem.\ Phys.\ 99, 2591 (1993).
[5] Spectra Technology, \emph{Experimental Support for KrF Laser
Modeling}, final report to Los Alamos, 19 July 1989, OSTI 10136715.
[6] Hakuta et al., Jpn.\ J.\ Appl.\ Phys.\ 29, 1194 (1990). [7] Burnham and Searles, J.\ Chem.\ Phys.\ 67, 5967 (1977).
[8] Morozov, Heindl, Kr\"ucken, Ulrich and Wieser, J.\ Appl.\ Phys.\ 103,
103301 (2008).\par}}}
\medskip

LIGHT is a 0-D model of krypton and fluorine under a power drive. Its basic
network has \ov{rates.nBasic} reaction rows. Every KrF rate among them but
one is a measurement of that process or a joint fit to kinetic measurements
of it. The exception is the KrF(B) radiative lifetime, where LIGHT keeps
Kannari's \lm{tauB}\unit{ns} and Eden measured 6.8 $\pm$ 0.2\unit{ns}. No
rate was set by a laser measurement. The one independent Kr/F$_2$
laser LIGHT has been scored against is Tahoma, and no fit used any number
from it. There LIGHT gives the small-signal gain at \x{\ov{tah.g0.x}} the
measured value, the saturation intensity at \x{\ov{tah.isat.x}} and the
formation efficiency \gIP{} at \x{\ov{tah.eta.x}}. The largest open
questions are three links of the formation chain that no one has measured
and the trimer rate above 600\unit{K}. Table~\ref{tab:score} lists every test.

\begin{table}[h]
\centering\footnotesize
\renewcommand{\arraystretch}{1.12}
\begin{tabularx}{\linewidth}{@{}>{\raggedright\arraybackslash}X
  >{\raggedright\arraybackslash}p{2.6cm}>{\raggedright\arraybackslash}p{2.5cm}
  >{\raggedright\arraybackslash}p{2.2cm}@{}}
\toprule
test & measured & LIGHT & role of the data\\
\midrule
energy ledger closure, \ov{closure.runs} runs across the modes & & below
  \ov{closure.max} of \Edep & bookkeeping\\
KrF(B) decay against Kr and F$_2$, Eden 1978 & & rms \ov{sc.eden_kr} and
  \ov{sc.eden_f2}\,\% & train\\
KrF(B,C) slow decay, Gadomski 1992 Figs. 1 and 2 & & rms \ov{sc.ksu_fig1}
  and \ov{sc.ksu_fig2}\,\% & train\\
C/B emission ratio at 300\unit{K}, Yu and Setser 1990 & \ov{ys.meas}
  $\pm$ \ov{ys.err} & \ov{sc.ys_ir300}\,\% & train\\
KrF(B,C) decays, Xu 1993 Figs. 5 to 7 & & rms \ov{sc.xu_figs5_7}\,\% &
  held out\\
decay constants Xu 1993 printed, 50 and 100 Torr Kr & \ov{xu.50.meas}
  and \ov{xu.100.meas}\unit{s^{-1}} & \x{\ov{xu.50}} and \x{\ov{xu.100}} &
  structure check\\
C-A decays, Hakuta 1990 & & rms \ov{sc.hakuta_fig2} and
  \ov{sc.hakuta_fig3}\,\% & held out, in conflict\\
$\tau_B$, Burnham and Searles 1977 & \ov{sc.bs.meas} $\pm$
  \ov{sc.bs.err}\unit{ns} & \ov{sc.bs_tau_b}\,\% & held out, in conflict\\
C/B ratio 230\unit{K} over 300\unit{K}, Yu and Setser & \ov{ysr.meas}
  $\pm$ \ov{ysr.err} & \ov{sc.ys_ratio_230_300}\,\% & held out\\
\midrule
\Krs{} + \Krp{} made per eV, pure Kr and Ar, Saito 2003 & & \x{\ov{pure.Kr.x}}
  and \x{\ov{pure.Ar.x}} & check\\
\Krtp{} + \el{} recombination, Shiu and Biondi 1977 & & \x{\ov{pure.dr.lo}}
  to \x{\ov{pure.dr.hi}} & check\\
\Krp{} + 2Kr \ra \Krtp{} + Kr, Smith et al. 1972 & \ov{pure.conv.lo} to
  \ov{pure.conv.hi} & \ov{pure.k15} & check\\
VUV yield of pure Kr at low pump, Morozov 2008 & \ov{pure.vuv.meas}
  $\pm$ \ov{pure.vuv.err} & \ov{pure.vuv}, \x{\ov{pure.vuv.x}} & check\\
electron closure at Tahoma, kinetic EEDF loop & $T_e$ \ov{el.loop.te}\unit{eV}
  & \ov{el.light.te}\unit{eV} & reference\\
\midrule
Tahoma small-signal gain $g_0$ & \ov{tah.g0.meas}\unit{\%/cm} &
  \ov{tah.g0}, \x{\ov{tah.g0.x}} & device, score only\\
Tahoma saturation intensity $I_\mathrm{sat}$ & \ov{tah.isat.meas} $\pm$
  \ov{tah.isat.err}\unit{MW/cm^2} & \ov{tah.isat}, \x{\ov{tah.isat.x}} &
  device, score only\\
Tahoma formation efficiency \gIP & \ov{tah.eta.meas} & \ov{tah.eta},
  \x{\ov{tah.eta.x}} & device, score only\\
\bottomrule
\end{tabularx}
\caption{LIGHT against the data. Rates are in cm$^3$s$^{-1}$ or
cm$^6$s$^{-1}$. A percentage is LIGHT's deviation from the data, signed
for a single value and rms for a set. A \x{} factor is LIGHT over the
measurement.}
\label{tab:score}
\end{table}

\section{How the numbers were made}

Each dataset has one role. Kinetic measurements of the channels being
fitted train them. Other kinetic measurements are held out and test the
fit. Laser performance only scores, and the fitting code refuses a dataset
of that kind. The train set is the largest group of kinetic datasets that
agree with one another within their errors, found by leaving each group out
in turn. This is the repository's predictive-only policy. Measurements
score the model after the fact and never steer a number. Every number in
this note is written by \texttt{docs/make\_overview.py} from LIGHT and from
\texttt{ML/}, and a test fails when those files no longer match the model.

\section{The KrF(B,C) block}

KrF has two emitting states. B lases at 248\unit{nm}, and C radiates slowly
in a band near 275\unit{nm}. LIGHT treats them as the two coupled levels of
eq. 5 of Gadomski et al. (1992) and reproduces that model's analytic decay
rate to within \ov{eq5.dev} at every pressure tested. The block has \ov{fit.n} constants.
Nine come from one joint fit to three measurement sets that agree with one
another, the KrF(B) decays of Eden et al. (1978), the Kansas State B/C
decays (Gadomski et al. 1992) and Yu and Setser's equilibrium C/B ratio
(1990). The tenth, $\tau_B$, stays at Kannari's \lm{tauB}\unit{ns}, inside
the fit's interval of \ov{fit.tau_B.lo} to \ov{fit.tau_B.hi}\unit{ns}. The fit's median and 68\,\% interval for B-C mixing by Kr is
$k_{BC}=$ \ov{fit.k_BC.mid} [\ov{fit.k_BC.lo}, \ov{fit.k_BC.hi}]
cm$^3$s$^{-1}$. The equilibrium ratio is $K_\mathrm{eq}=$
\ov{fit.K_eq.mid} [\ov{fit.K_eq.lo}, \ov{fit.K_eq.hi}], and the C lifetime is
$\tau_C=$ \ov{fit.tau_C.mid} [\ov{fit.tau_C.lo}, \ov{fit.tau_C.hi}]\unit{ns}.
LIGHT carries the best fit, $k_{BC}=\lm{k.BC}$\unit{cm^3s^{-1}},
$K_\mathrm{eq}=$ \lm{Keq} and $\tau_C=$ \lm{tauC}\unit{ns}. Mixing is slow, so KrF made in C mostly stays there. At 1\unit{atm} and
300\unit{K} the lasing B state holds \ov{fit.bshare.mid}
[\ov{fit.bshare.lo}, \ov{fit.bshare.hi}] of the manifold, and LIGHT's
rounded constants give \ov{fit.bshare.light}.

The fitted decay sets are matched to between \ov{sc.train.lo} and
\ov{sc.train.hi}\,\% rms. The held-out decays of Xu et al. (1993) are
matched to \ov{sc.xu_figs5_7}\,\%, which is inside the 30\,\% by which the
Kansas State figures scatter at one condition. The two decay constants Xu
et al. print come out at \x{\ov{xu.50}} and \x{\ov{xu.100}}, where the notes'
list as written, with no C state, gives \x{\ov{xu.50.v0}} and
\x{\ov{xu.100.v0}}. Three held-out measurements disagree. Hakuta et al.
(1990) measured faster C-A decays, missed by \ov{sc.hakuta_fig2} and
\ov{sc.hakuta_fig3}\,\% rms. Their data also conflict with the three fitted
sets, which is why they are held out. Burnham and Searles (1977) measured
$\tau_B=$ \ov{sc.bs.meas}\unit{ns}, where LIGHT uses Kannari's
\lm{tauB}\unit{ns} and Eden measured 6.8\unit{ns}. Yu and Setser's
C/B ratio between 230 and 300\unit{K} comes out \ov{sc.ys_ratio_230_300}\,\%,
because the energy ledger books a B-C gap of \lm{gapBC}\unit{cm^{-1}} where
their two ratios imply 180 $\pm$ 30\unit{cm^{-1}}.

Tahoma's gain is the sharpest test of the block, since $g_0$ is the
KrF(B) density times the stimulated cross section. Over \ov{post.n} draws from the
fit's posterior the gain is \ov{post.g0.mid}\unit{\%/cm} [\ov{post.g0.lo},
\ov{post.g0.hi}] against the measured \ov{tah.g0.meas}.

\section{KrF formation}

The chain from the beam's \Krp{} and \Krs{} to KrF sets the formation
efficiency. LIGHT carries every link that has been measured at its
measured value. \Krs{} + \Ft{} reacts at $\lm{k.harp}$\unit{cm^3s^{-1}}, the
mean of three first-hand measurements. Velazco, Kolts and Setser (1976)
measured the KrF$^*$ yield of that reaction as \lm{pc.harpY}\,\%, the
KrF$^*$ formation rate constant over the total quenching, and published it
as unity $\pm$ 0.20. Yu and Setser (1990) put \lm{pc.fC}\,\% of that
KrF$^*$ in C. The \Krts{} route carries Ueda's \Krs{} + 2Kr rate, the
measured 270\unit{ns} \Krts{} lifetime and Quigley and Hughes' \Krts{} +
\Ft{} rate. Appendix A gives every value with its source.

Tahoma measured \gIP{} $=$ \ov{tah.eta.meas}. ML puts a 248-nm probe flux
on LIGHT and fits its $I_\mathrm{sat}$ at fixed $g_0$, as Tahoma did. That
gives \ov{tah.eta}. Its KrF formation rate in 5-eV quanta is
\ov{tah.form} of the pump. The difference is KrF that forms in C, or mixes
into it, and decays there without meeting the laser field.

Table~\ref{tab:links} moves one link at a time to another value it could
take. The measured range of the harpoon yield, 0.70 to 1.0, moves \gIP{}
from \ov{link.yield.lo} to \ov{link.yield.hi} and brackets Tahoma's
\ov{tah.eta.meas}. Three links have no measurement at all. They are the
KrF$^*$ yield of ion-ion recombination with F$^-$, the share of that KrF
formed in C, and \Krts{} + F \ra KrF. Only analogs exist for the first
two, and both analogs lower \gIP{} toward the measurement. Removing the third
lowers it too. The predictive-only
policy lets no analog set a link and no link be chosen by its fit to
Tahoma. The rest of the \Krts{} route moves \gIP{} by \ov{link.minor} at
most.

\begin{table}[H]
\centering\footnotesize
\renewcommand{\arraystretch}{1.12}
\begin{tabularx}{\linewidth}{@{}>{\raggedright\arraybackslash}X
  >{\raggedright\arraybackslash}p{1.8cm}>{\raggedright\arraybackslash}p{1.9cm}
  >{\raggedright\arraybackslash}p{3.3cm}r r@{}}
\toprule
link & LIGHT & tried & what the value is & \gIP & change\\
\midrule
\linksTable
\bottomrule
\end{tabularx}
\caption{Each formation link moved alone from LIGHT's chain at Tahoma,
\ov{tah.f2torr}\unit{Torr} of F$_2$ in Kr at \ov{tah.torr}\unit{Torr} in all,
pumped at \ov{tah.pin}\unit{kW\,cm^{-3}}. LIGHT's chain gives \ov{link.base}, and
Tahoma measured \ov{tah.eta.meas}.}
\label{tab:links}
\end{table}

\section{The electron side}

LIGHT integrates each krypton electron-impact rate over an EEDF with two
temperatures. The bulk is a Maxwellian at $T_e$. Above the 9.915-eV threshold the
tail has its own temperature, read from a table of kinetic EEDFs with
electron-electron collisions at e-beam-sustained states. Each electron that
is captured takes the energy it really carries, $\langle\varepsilon\sigma
v\rangle/\langle\sigma v\rangle$. For \Ft{} attachment at Tahoma that is
\ov{el.eatt}\unit{eV}, where $\tfrac32 T_e$ would book \ov{el.eatt.32}\unit{eV}.
A kinetic loop that iterates LIGHT with a Boltzmann solution is the
reference. At Tahoma LIGHT's own closure lands within \ov{el.dev.te}\,\% of
the loop's $T_e$ and \ov{el.dev.ne}\,\% of its $n_e$, and within
\ov{el.dev.g0} and \ov{el.dev.formation}\,\% of its gain and formation.

The cross sections were checked against swarm data. With Biagi's krypton
set the computed $W$ is \ov{sw.krw.lo} to \ov{sw.krw.hi} times Nakamura's
measured values. For \Ft{} attachment the rate constants measured by
different groups disagree beyond their errors, and the SIGLO set LIGHT
uses gives between \ov{sw.f2.lo} and \ov{sw.f2.hi} times the measured
values. The two
groups that agree give a scale of \ov{sw.slow} $\pm$ \ov{sw.slow.err} on the
main resonance, \ov{sw.slow.sigma} standard deviations from SIGLO. LIGHT
keeps SIGLO as it is.

On pure gases the fluorine-free half of the chain checks out. The beam
drive's Spencer-Fano yields make \ov{pure.Kr.x} times the \Krs{} and \Krp{}
per eV that the measured $W$ and $N_\mathrm{ex}/N_i$ give in Kr, and
\ov{pure.Ar.x} times in Ar. LIGHT's \Krtp{} recombination is
\ov{pure.dr.lo} to \ov{pure.dr.hi} times Shiu and Biondi's between
\ov{pure.dr.te.lo} and \ov{pure.dr.te.hi}\unit{eV}. The VUV yield of pure Kr at low pump is the one
miss. LIGHT gives \ov{pure.vuv} where Morozov et al. (2008) measured
\ov{pure.vuv.meas} $\pm$ \ov{pure.vuv.err}. Without the thermal electrons'
excitation LIGHT gives \ov{pure.vuv.nothermal}, close to the one photon per
excitation and per ion that the measured $W$ and $N_\mathrm{ex}/N_i$ allow,
\ov{pure.vuv.ideal}. No rate in LIGHT removes the rest. The measured e-beam
yields of Ne, Ar, Kr and Xe all fall below that ideal, so the shortfall is
either physics LIGHT lacks or a systematic error of the measurement. In
pure Kr at low pump the bulk EEDF is not Maxwellian either, and the
one-temperature energy equation cannot follow it there. That limits the
pure-gas check and not KrF.

\section{The energy ladder}

Figure~\ref{fig:ladder} places every species of v0 at its formation
energy, the energy above ground-state Kr and \Ft{} at rest that
Appendix~A lists, and draws the rows that move energy between them. It
follows the ladder view of the website's excimer simulation, with v0's own
numbers. v0 keeps REBUILD's estimated KrF(B), \lad{E.B}\unit{eV}, which
puts it \lad{DeB}\unit{eV} below the ion pair \Krp{} + \Fm{} and
\lad{BoverKrF}\unit{eV} above Kr + F. The \lad{wl.G4}-nm photon carries
\lad{hv.G4}\unit{eV} of that and leaves KrF(X) \lad{heat.G4}\unit{eV}
above Kr + F, which heats the gas as X falls apart. The harpoon G3 heats
the gas by \lad{heat.G3}\unit{eV} for each KrF(B) it makes and the
trimer's formation by \lad{heat.G5}\unit{eV}. \KtF{} emits
\lad{hv.C3}\unit{eV} at \lad{wl.C3}\unit{nm} and leaves
\lad{heat.C3}\unit{eV}. The \Krts{} route, most of the charge chemistry
and the \Ft{} electron rows are optional groups in v0, drawn dashed, and
the species only they make are pale.

\begin{figure}[p]
\centering
\ladpicture
\caption{The energy ladder of LIGHT v0, version \lad{version}, in eV above
ground-state Kr and \Ft{} at rest. A bar is a species at its formation
energy and a dashed bar a pair, the sum of two. An arrow is a row,
labelled with its id and solid in the default network. The rows of the
optional groups \lad{groups} are dashed, and the species only they make,
\lad{pale}, are pale. A photon's arrow drops by its quantum, and the red
squiggle below it is the heat its lower state gives the gas as it falls
apart. A red stub is a quench. The energy a row releases heats the gas,
except a photon's, which leaves, and that of \lad{toe}, which goes to the
electrons. KrF(B) and KrF(C) lie \lad{gapBC}\unit{cm^{-1}} apart and are
drawn as a pair of bars that every B and C row meets. The rows
\lad{notdrawn} move nothing between levels and are not drawn.}
\label{fig:ladder}
\end{figure}

\section{The system map}

The map on the next page draws the basic network's energy flow at the
notes' reference point, \lm{p}\unit{atm} of Kr with \lm{xF2}\,\% F$_2$,
\lm{P}\unit{kW\,cm^{-3}} for \lm{pulse}\unit{\mu s} and
$n_e=\lm{ne}$\unit{cm^{-3}} held. Band widths are proportional to the
energy they carry, and every box conserves it. Of the deposited energy,
\lm{krf.eff}\,\% passes through KrF in 5-eV quanta. Of that KrF,
\lm{brf.G4}\,\% radiates at 248\unit{nm}, \lm{brf.tri}\,\% forms the trimer
and \lm{brf.e}\,\% is quenched by electrons, so \lm{out.248}\,\% of the
deposited energy leaves at 248\unit{nm}. The notes' six reactions as
written, drawn in \texttt{docs/system\_map.pdf}, give \lm{om.v0}\,\%. The
difference is the C state and the collisional losses that the KrF
measurements require.

% the ladder is a float, which pdfpages would leave pending past the map
\clearpage
\includepdf[pages=1,fitpaper=true]{system_map_basic.pdf}

\section{Where LIGHT is not yet reliable}

\begin{itemize}
\item The trimer. KrF + 2Kr \ra \KtF{} takes \lm{brf.tri}\,\% of the KrF
at the reference point.
Its $(T_g/300)^{\lm{tpow.G5}}$ law is Shui's calculation for 300 to
600\unit{K}, measured only between 315 and 420\unit{K}. At the notes'
reference point the gas ends at \lm{TgEnd}\unit{K}.
\item The drive. The notes' power drive feeds the electron gas directly and
borrows the tail table made at e-beam states. LIGHT has no e-beam route of
its own, and the Tahoma scores ran under a prototype beam drive that is
not part of LIGHT's model.
\item Fluorine. At the reference point \Ft{} falls below 10\,\% of the fill
at \lm{F2burn}\unit{ns}, inside the pulse, and LIGHT has no F + F
recombination.
\item Ions. With $n_e$ held fixed and the ions group off, nothing
recombines, so the ion store is bookkeeping.
\item The C-A band centre near 275\unit{nm} is read from a figure to
$\pm$10\unit{nm}.
\item Hakuta's C-A decays and Burnham and Searles' $\tau_B$ remain
unexplained conflicts.
\end{itemize}

\appendix
\section{Reaction rates}

Rates are the values at 300\unit{K}. The column $T$ gives the law LIGHT
applies to them. Electron-impact rows use $\langle\sigma v\rangle$ over the
EEDF, and Section A.3 lists their Maxwellian values. Each row's
full citation is in \texttt{v0/network.py}. The row ids are the ones the
code and the system map use.

\footnotesize
\setlength{\tabcolsep}{3.5pt}
\renewcommand{\arraystretch}{1.12}

\subsection{The basic network, \ov{rates.nBasic} rows}
\begin{longtable}{@{}l>{\raggedright}p{4.1cm}>{\raggedright}p{3.1cm}
  >{\raggedright}p{1.9cm}>{\raggedright\arraybackslash}p{4.75cm}@{}}
\toprule
row & reaction & rate & $T$ & source\\
\midrule
\endhead
\bottomrule
\endfoot
\ratesBasic
\end{longtable}

\subsection{Optional groups on the basic network, \ov{rates.nGroups} rows}
The groups \ov{rates.groups} are off at the reference point and switch on
with \texttt{-{}-add}. The first four \Krts{} rows carry the basic network's
measured values.
\begin{longtable}{@{}l>{\raggedright}p{4.1cm}>{\raggedright}p{3.1cm}
  >{\raggedright}p{1.9cm}>{\raggedright\arraybackslash}p{4.75cm}@{}}
\toprule
row & reaction & rate & $T$ & source\\
\midrule
\endhead
\bottomrule
\endfoot
\ratesGroups
\end{longtable}

\subsection{Electron-impact rate coefficients}
Maxwellian $\langle\sigma v\rangle$ in cm$^3$s$^{-1}$ at the $T_e$ of each
column. The last column is the factor the two-temperature tail applies at
the reference point's mid-pulse state, $T_e=$ \ov{rates.teMid}\unit{eV} with
the tail at \ov{rates.rMid}\,$T_e$. The rows below the rule give the energy
an attached electron takes from the electron gas.
\begin{center}\scriptsize
\begin{tabular}{@{}llrrrrrr@{}}
\toprule
row & process & \multicolumn{5}{c}{$T_e$ (eV)} & tail\\
\cmidrule(lr){3-7}
 & & \ratesTe & factor\\
\midrule
\ratesXsec
\bottomrule
\end{tabular}
\end{center}

\subsection{Where v0, the notes' list as written, differs}
The v0 network keeps the notes' six reactions with REBUILD's values, and its
optional groups, \ov{rates.v0groups}, are REBUILD's. These rows differ from
the basic network or exist only in v0. The rows ending in \_C are the bc
group's copies for KrF(C).
\begin{longtable}{@{}l>{\raggedright}p{4.1cm}>{\raggedright}p{3.1cm}
  >{\raggedright}p{1.9cm}>{\raggedright\arraybackslash}p{4.75cm}@{}}
\toprule
row & reaction & rate & $T$ & source\\
\midrule
\endhead
\bottomrule
\endfoot
\ratesVzero
\end{longtable}

\subsection{Formation energies}
Energies above ground-state Kr and F$_2$ at rest. Every reaction's energy
release is the difference of these numbers, so no row can make or lose
energy. $c_v$ is the heat capacity per particle in units of $k_B$.
\begin{center}
\begin{tabular}{@{}lrrl@{}}
\toprule
species & $E_\mathrm{form}$ (eV) & $c_v$ & what sets it\\
\midrule
\speciesTable
\bottomrule
\end{tabular}
\end{center}

\vfill
{\color{ink2} Generated from LIGHT v0, version \lm{version}, and ML by
\texttt{docs/make\_overview.py} and \texttt{docs/make\_system\_map\_basic.py}.
No number in this document is typed by hand.}

\end{document}
